\scanpage{11}
% archive foliation: 10
\[
\rho(n)\overset{\mathrm{df}}=\operatorname{rg}_kL_n
=\sum_{1\leq i\leq n+1}\mu_i\binom{n+1}{i}
+\sum_{0\leq j\leq n}\nu_j\binom{n}{j}
\]
\diagram{\begin{tikzcd}[column sep=huge]
\binom{n+1}{i}\arrow[d] & \binom{n}{j}\arrow[d]\\
\text{pol. de degré $i$ en $n$} & \text{pol. de degré $j$ en $n$}
\end{tikzcd}}
où $\mu_i=\operatorname{rg}_kM_i$, $\nu_j=\operatorname{rg}_kN_j$.

donc si \deleted{\unsure{$\mu_i\ (\nu_j)$}} un $\mu_i$ ou un $\nu_j$
est $\neq0$, cela montre que $\rho(n)$ pour
$n$ grand est minoré par un polynôme
de degré $i$ resp $j$.

[2°) Cas de $\Lambda^k$. Utilisant la formule
du $\Lambda^k$ d'une somme directe et
1°), on est ramené à prouver que
le $\Lambda^k$ d'un module $\frac12$ simplicial
$L_*$ de la forme $\Lambda^i\Phi_*\otimes M$, $\Lambda^j\Psi_*\otimes N$.

Comme $M$, $N$ sont facteurs directs d'un
libre $P$, \inserted{$L_*$ est facteur direct de
$L'_*=\Lambda^i\Phi_*\otimes P$ \unsure{ou} $\Lambda^j\Psi_*\otimes P$,}
\deleted{\unsure{et}} \deleted{\unsure{cela montre}}
\inserted{\unsure{la propriété du cor. 1 et du cor 2}
\deleted{\unsure{prop}} est stable par fact. direct}
\deleted{\unsure{la propriété de déf.}}
et en utilisant la
formule du $\Lambda^k$ d'une somme directe,
\unsure{montre} que $\Lambda^k(L_*)$ est facteur
direct de $\Lambda^kL'_*$. On est ramené à
prouver $\Lambda^kL'_*$ est \deleted{\unsure{str.}} spécial, donc
\unsure{supp.} $M$, $N$ libres --- donc $=$ (par la
formule de la somme) $M=k$, $N=k$.
Donc on est ramené à prouver
